\section{Introduction}
\label{sec:intro}

RNG~\cite{rng} replaces the layered fat tree with a single layer of top-of-rack (ToR) routers wired as a random graph. Each router has $d$ links to other routers ($d = 64$ in the paper), and a 1{,}000-router fabric connects 64{,}000 servers with 45\% fewer switches than a comparable fat tree. Its routing protocol, Spraypoint, makes such a graph usable on commodity switches: the first router \emph{sprays} each flow to a random neighbor, and from there every router \emph{points} it toward the destination through a set of rings. This gives each pair of racks dozens of independent paths and spreads load across the whole fabric.

Spraying is tuned for throughput. Much datacenter traffic, though, is small and latency-sensitive: RPCs, key-value lookups, control messages. For these, what matters is the tail of the flow completion time (FCT), not bandwidth, and spraying has a cost. A message between two directly connected racks reaches the destination's neighbor only once in $d$ sprays, so it usually takes three or four hops instead of one. Every extra hop is another queue to wait in and another place to be dropped.

This project asks whether a simple change helps: let small flows skip the spray and take the shortest path. The paper itself warns against shortest paths between neighbors, because they concentrate traffic on direct links. So the question is a trade-off:

\begin{quote}
\emph{On an RNG fabric, does routing small flows on shortest paths, instead of spraying them, reduce their tail latency? What does it cost the rest of the traffic, and when does the cost outweigh the gain?}
\end{quote}

Our proposal stated four hypotheses:
\begin{description}[leftmargin=2.5em,style=nextline]
  \item[H1 (latency)] For small flows between neighboring routers, the bypass cuts the FCT tail, because they cross one queue instead of two to five.
  \item[H2 (no free lunch)] At low small-flow volume, other traffic barely notices. As the small flows' share of bytes or the load grows, direct links become hotspots, and past a crossover the bypass makes things worse.
  \item[H3 (reach)] Only about $d/(n-1)$ of router pairs are neighbors (6.4\% at the paper's scale), so to reach most small flows the bypass must cover every pair, on shortest paths. That raises the hotspot risk.
  \item[H4 (overhead)] The bypass costs a switch one extra table lookup, negligible next to queueing.
\end{description}

\paragraph{Where each hypothesis is tested.}
\begin{itemize}
  \item \textbf{H1}: small-flow latency under uniform traffic (\secref{sec:h1}).
  \item \textbf{H2}: the cost to other traffic, under uniform traffic (\secref{sec:h2-uniform}) and under worst-case traffic (\secref{sec:h2-perm}).
  \item \textbf{H3}: small flows between non-neighbors (\secref{sec:h3}).
  \item \textbf{Models}: the paper's analytical models, extended to the bypass, against the simulation (\secref{sec:models-vs-sim}).
  \item \textbf{H4}: argued from the switch pipeline, not measured (\secref{sec:h4}).
\end{itemize}

\paragraph{Roadmap.}
\secref{sec:background} explains RNG and Spraypoint for a reader new to them, and \secref{sec:bypass} defines the bypass. \secref{sec:method} describes how we measure: the simulator, the traffic, the metrics, and a flow-level checker. \secref{sec:models} extends the paper's models to the bypass. \secref{sec:results} presents the results in the order the project found them, \secref{sec:discussion} turns them into a policy, and \secref{sec:conclusion} concludes. The appendices hold the full model derivations, the validation of our Spraypoint implementation, and how to reproduce every number.
